\section{Finite relational viability and coherent witness clouds}
\label{sec:relations}

Nonunique witnesses often appear at every finite stage.  Choosing one witness
independently at each stage may produce an incompatible sequence even when a
compatible finite run exists.  DARC resolves the finite problem by backward
viability before making any choice.

\subsection{Layered finite relations}

Fix a horizon \(H\).  For each \(0\le h\le H\), let
\(X_h=\{0,\ldots,B_h-1\}\) be a finite coded set, and for
\(h<H\) let \(E_h(x,y)\) be a reflected compatibility relation between
\(X_h\) and \(X_{h+1}\).  Let \(G_H(x)\) be the terminal acceptance bit.
Define viability backward:
\begin{align*}
  V_H(x)&=G_H(x),\\
  V_h(x)&=\bigvee_{y<B_{h+1}}
      \bigl(E_h(x,y)\wedge V_{h+1}(y)\bigr).
  \tag{11.1}\label{eq:viability}
\end{align*}
The finite disjunction is compiled to a primitive-recursive bit.  Since the
horizon is fixed or input-bounded with a fixed step definition, the family can
be generated by primitive recursion on the remaining horizon.

\begin{theorem}[Finite viability]
\label{thm:finite-viability}
For each \(h\le H\) and \(x\in X_h\),
\[
  V_h(x)=1
  \quad\leftrightarrow\quad
  \text{there exist }x_{h+1},\ldots,x_H
  \text{ forming an }E\text{-compatible run ending in }G_H.
\]
\end{theorem}

\begin{proof}
Backward induction on \(H-h\).  At the terminal layer, the claim is the
reflection theorem for \(G_H\).  At layer \(h\), the finite-disjunction
reflection theorem turns \(V_h(x)=1\) into existence of a
\(y<B_{h+1}\) with \(E_h(x,y)\) and \(V_{h+1}(y)=1\).  Apply the induction
hypothesis to append a compatible tail.  Conversely, the first element of any
compatible tail provides a true disjunct, and the induction hypothesis proves
its viability bit.
\end{proof}

\subsection{Least compatible path}

Once viability is known, define the successor at layer \(h\) to be the least
\(y<B_{h+1}\) satisfying
\[
  E_h(x,y)\wedge V_{h+1}(y).
\]
This is a bounded least search.  Iterating the successor for \(H-h\) steps
gives a canonical finite path.

\begin{theorem}[Canonical finite path]
\label{thm:least-path}
If \(V_0(x_0)=1\), the least-successor construction yields a coded path
\(x_0,\ldots,x_H\) satisfying every edge relation and the terminal goal.
Each selected successor is least among viable successors.
\end{theorem}

\begin{proof}
By \cref{thm:finite-viability}, a viable successor exists at every nonterminal
state reached from a viable state.  The bounded-search theorem selects its
least code and proves validity.  Equation \eqref{eq:viability} ensures the
selected successor is itself viable.  Quantifier-free induction on the number
of steps proves every transition and terminal acceptance.
\end{proof}

The path may remain a local certificate.  It is named only if the public query
requests it.

\subsection{Backward viability before choice}

Consider a relation in which each layer has many individually valid elements
but only a few extend to the terminal goal.  A forward greedy choice based only
on local validity can get stuck.  Backward viability removes all dead ends
first.  Therefore the subsequent least choice is compatible by construction.
This resolves the ``independent finite witnesses'' difficulty without any
appeal to dependent choice.

\subsection{Stable observations on witness clouds}

Often the path itself is irrelevant.  Suppose each state has an observation
\(O_h(x)\), and compatible edges satisfy one of the following:
\begin{align*}
  O_{h+1}(y)&=O_h(x), &&\text{exact stability},\\
  O_{h+1}(y)&\sim O_h(x), &&\text{quotient stability},\\
  d(O_{h+1}(y),O_h(x))&\le\varepsilon_h,
     &&\text{metric stability}.
\end{align*}
Then any viable finite path yields a root observation.  In the metric case,
triangle inequality gives
\[
  d(O_H(x_H),O_0(x_0))
  \le \sum_{h<H}\varepsilon_h.
\]
This conclusion is finite and constructive.  It avoids selecting an
infinitary quotient representative.

\begin{theorem}[Finite observational descent]
\label{thm:observational-descent}
Assume \(V_0(x_0)=1\) and edge-wise exact, quotient, or metric compatibility.
Then every viable terminal witness has an observation related to
\(O_0(x_0)\) by the corresponding composed relation; in the metric case the
explicit error is the finite sum of edge errors.
\end{theorem}

\begin{proof}
Use \cref{thm:least-path} or keep an arbitrary viable path scoped.  Induct along
the finite path.  Exact equality composes by transitivity; a declared
quotient relation must be supplied with transitivity; metric errors compose by
triangle inequality and finite summation.  The certificate can then be sealed
if the conclusion mentions only the observation.
\end{proof}

\subsection{The infinite-path barrier}

For every \(H\), existence of a compatible path of length \(H\) does not by
itself give a primitive-recursive infinite path in \(\iPRA\).  The least path
for horizon \(H+1\) need not extend the least path for horizon \(H\).  Even a
coherent family of nonempty finite levels requires an extension principle or a
primitive-recursive section to choose one successor forever.

DARC therefore recognizes three legitimate endpoints:
\begin{enumerate}
  \item a theorem for each finite horizon;
  \item a coherent path generated by an explicit primitive-recursive extension
        function, with its totality proof;
  \item a conditional theorem with the extension producer left residual.
\end{enumerate}
It rejects the unqualified leap from horizonwise existence to an infinite
object.

\subsection{The regular-name barrier}

Even a coherent infinite path of approximants does not automatically define a
constructive real name.  One also needs a quantitative tail condition, for
example
\[
  m,n\ge M(p)\Longrightarrow d(x_m,x_n)\le 2^{-p},
\]
with primitive-recursive modulus \(M\).  DARC represents \(M\) by a precision
transfer or residual producer.  Without it, finite observational descent
remains valid, but a global Cauchy-name theorem is not emitted.

\subsection{Applications}

The same pattern applies to:
\begin{itemize}
  \item compatible local coordinate charts;
  \item branch choices in algebraic continuation over a finite mesh;
  \item finite tiling patches and boundary states;
  \item dynamic-programming witnesses;
  \item finite model approximations;
  \item staged Diophantine certificates whose constraints must agree across
        projections.
\end{itemize}
In each case, viability separates finite compatibility from unjustified
infinite choice.
